
\subsubsection{RQ1: Batch Normalization Amplifies Worst-Group Gaps}

At 90\% average accuracy, ResNet-18-BN shows mean worst-group gap of 19.92 $\pm$ 2.18 pp across 10 seeds, while ResNet-18-LN shows 10.51 $\pm$ 2.58 pp. Gap difference: 9.41 pp.

\begin{table}[htbp]
\centering
\resizebox{\columnwidth}{!}{%
\begin{tabular}{llll}
\toprule
Architecture & Mean Gap (pp) & Std Dev (pp) & Seeds Reaching 90\% \\
\midrule
ResNet-18-BN & 19.92 & 2.18 & 10/10 \\
ResNet-18-LN & 10.51 & 2.58 & 10/10 \\
Difference & 9.41 & --- & --- \\
\bottomrule
\end{tabular}
}
\end{table}

Statistical test: t(9) = 7.14, p = 5.$43\times 10^{-5}$, Cohen's d = 3.94

All 10 seeds for both architectures reached 90\% average accuracy. The null hypothesis (gap difference < 5.0 pp) is rejected (p < 0.001). Effect size (d = 3.94) exceeds threshold (d $\geq$ 0.8).

\subsubsection{RQ2: Gradient Mechanism}

During epochs 0--19, ResNet-18-BN exhibits mean gradient ratio 1.2032 $\pm$ 0.0808, while ResNet-18-LN shows 0.9532 $\pm$ 0.0579. Difference: 0.2500 (26.23\% increase).

\begin{table}[htbp]
\centering
\resizebox{\columnwidth}{!}{%
\begin{tabular}{lll}
\toprule
Architecture & Mean Gradient Ratio & Std Dev \\
\midrule
ResNet-18-BN & 1.2032 & 0.0808 \\
ResNet-18-LN & 0.9532 & 0.0579 \\
\% Increase & 26.23\% & --- \\
\bottomrule
\end{tabular}
}
\end{table}

Statistical test: t(9) = 9.66, p < 0.001, Cohen's d = 4.32

Batch Normalization shows 26.23\% higher gradient flow to spurious-aligned samples during early training. This provides mechanistic evidence for why Batch Normalization amplifies worst-group gaps.

\subsubsection{RQ3: Ranking Consistency}

On synthetic datasets, architecture ranking shows Spearman $\rho$ = 1.0000 with zero rank reversals. Both synthetic Waterbirds and mock CelebA rank Layer Normalization (lower gap) ahead of Batch Normalization.

\begin{table}[htbp]
\centering
\resizebox{\columnwidth}{!}{%
\begin{tabular}{llll}
\toprule
Dataset & BN Gap (pp) & LN Gap (pp) & Ranking \\
\midrule
Synthetic Waterbirds & 19.92 & 10.51 & LN < BN \\
Mock CelebA & 17.96 & 9.20 & LN < BN \\
\bottomrule
\end{tabular}
}
\end{table}

Statistical test: $\rho$ = 1.0000, p-value not computable (n=2 architectures insufficient)

This result demonstrates that the ranking pipeline functions correctly. Scientific validation requires real CelebA training and at least 4 architectures. The planned attention mechanism hypothesis (h-m2) was incomplete due to technical failures, preventing inclusion of CBAM and ViT architectures.

\subsubsection{Effect Size Magnitude}

Gap difference (9.41 pp) exceeds planned threshold (5.0 pp) by 88\%. Effect sizes (d = 3.94 for RQ1, d = 4.32 for RQ2) are 4--$5\times$ larger than thresholds. Possible explanations:

\begin{enumerate}
\item Synthetic data (90\% spurious correlation) may amplify effects compared to real Waterbirds (85\%).
\item Constant learning rate may exaggerate BN-LN differences compared to scheduled learning rates.
\item Batch size 64 may interact with normalization type.
\end{enumerate}

Real dataset validation is required to determine actual effect magnitude.
